\paragraph{Implemented configuration}
{\color{blue}MGAT has $F=8$ input features per stock and day, three graph-attention layers of widths 16, 32, and 64 ($C=64$), and one encoder block (Eqs.~\eqref{eq:qkv}--\eqref{eq:ln2}) applied separately to each stock along the $T=20$ days of the window, with learned positional embeddings, four attention heads, a feedforward width of 128, and dropout 0.1. A linear output layer maps the last-day representation (Eq.~\eqref{eq:last_step_re}) of each stock to six outputs, the high, low, and closing log ratios for $t+1$ and $t+2$; the model has 44{,}846 trainable parameters (GCN and NCS: 44{,}622; MGNN: 44{,}850; graph-free Transformer: 2{,}718). Inputs and targets are standardised with the means and standard deviations of the training rows of each fold, and standardised inputs are clipped to $[-5,5]$. Temporal self-attention operates within each stock, and cross-stock information enters only through the graph-attention representation, which acts as inter-stock attention restricted to the relation edges.\par}

\subsection{Ranking-aware training objective}\label{subsec:objective}
{\color{blue}A forecast made at the close of day $t$ can be acted upon at the close of day $t+1$ at the earliest (Section~\ref{subsec:env}). The first return that such a decision earns is
\begin{equation}
q_{t,i}=\log\frac{C_{t+2,i}}{C_{t+1,i}}=y^{c}_{t,i,2}-y^{c}_{t,i,1},\qquad \hat q_{t,i}=\hat y^{c}_{t,i,2}-\hat y^{c}_{t,i,1},
\label{eq:q}
\end{equation}
where $y^{c}_{t,i,k}=\log(C_{t+k,i}/C_{t,i})$. For allocation across stocks, the cross-sectional ordering of $\hat q_{t,i}$ matters more than its level. The squared-error loss constrains $\hat q$ only indirectly, as the difference of two closing-price outputs whose errors are much larger than the variation of $q$ across stocks. MGAT is therefore trained with
\begin{equation}
\mathcal{L}=\frac{1}{|\mathcal{M}|}\sum_{(t,i,c)\in\mathcal{M}}\big(\hat y_{t,i,c}-y_{t,i,c}\big)^2+\lambda\Big(1-\frac{1}{|\mathcal{T}|}\sum_{t\in\mathcal{T}}\rho_t(\hat{\mathbf{q}}_t,\mathbf{q}_t)\Big),
\label{eq:loss}
\end{equation}
where the first term is the mean squared error over the six standardised targets on the set $\mathcal{M}$ of observed stock-days, $\rho_t$ is the Pearson correlation across the stocks with observed $q_{t,i}$ on day $t$, $\mathcal{T}$ contains the days in the mini-batch with at least five such stocks, and $\lambda=1$. The value of $\lambda$ was fixed before training and was not tuned. All neural forecasters in this study, including the comparators of Section~\ref{subsec3}, are trained with this objective. The linear comparator is a ridge regression of $q$ on the flattened 20-day window of the eight features, pooled over stocks, with the penalty chosen by the validation rank correlation of $q$; a second ridge regression on the same window forecasts the log range $\log(H_{t+1}/L_{t+1})$, with the penalty chosen by the validation mean squared error. Optimisation uses Adam with learning rate $10^{-3}$, batch size 32, at most 100 epochs, and early stopping with patience 10 on the validation loss. Each forecaster is trained with seeds 1--3, and the mean of the three forecasts is used downstream.\par}

\paragraph{Selection of the forecaster}
{\color{blue}The graph encoder was chosen before the evaluation year was examined, on the rank information coefficient of $\hat q$ (IC$_2$, the mean daily Spearman correlation between $\hat q_{t,i}$ and $q_{t,i}$) in the development years 2022 and 2023 (Section~\ref{subsec:splits}). Among the encoders that propagate information across stocks, graph attention had the highest development IC$_2$ in both sectors (banking 0.057, $t=3.5$; healthcare 0.018, $t=1.4$), compared with 0.054 and 0.010 for GCN and $-0.011$ and 0.012 for MGNN. The graph-free Transformer had a higher development IC$_2$ in banking (0.070), and it is therefore retained as a registered comparator, so that any contribution of the relation graphs is tested rather than assumed (Section~\ref{subsec:hypotheses}).\par}
